\section{Problemas}

\question{
Classifique e resolva (quando possível) cada um dos sistemas abaixo:
}

\subquestion{
\[
\begin{cases}
x + 2y + z = 0 \\
2x + 3y + z = 5 \\
x + y + 2z = 4
\end{cases}
\]
}
\answer{
\begin{align*}
\textbf{Classificação:}\ SPD
\end{align*}
\[
(x,y,z)=\left(\frac{19}{2},-\frac{9}{2},-\frac{1}{2}\right)
\]
}

\subquestion{
\[
\begin{cases}
x + y - z = -1 \\
2x + y - z = 1 \\
3x + y + 2z = 5 \\
4x + 2y + z = 4
\end{cases}
\]
}
\answer{
\begin{align*}
\textbf{Classificação:}\ SPD
\end{align*}
\[
(x,y,z)=\left(2,-\frac{7}{3},\frac{2}{3}\right)
\]
}

\subquestion{
\[
\begin{cases}
x + 3y - z = -1 \\
2x - y - 2z = -2 \\
x - 2y + z = 3
\end{cases}
\]
}
\answer{
\begin{align*}
\textbf{Classificação:}\ SPD
\end{align*}
\[
(x,y,z)=(1,0,2)
\]
}

\subquestion{
\[
\begin{cases}
x - y + z + t = 4 \\
2x - y - z = -3 \\
x - 2y + t = 1 \\
5x + z - t = 5
\end{cases}
\]
}
\answer{
\begin{align*}
\textbf{Classificação:}\ SPD
\end{align*}
\[
(x,y,z,t)=(0,-1,4,-1)
\]
}

\subquestion{
\[
\begin{cases}
x + 2y = 4 \\
2x + y + z = 3 \\
x - 2y - 3z = 0 \\
x + z = 3
\end{cases}
\]
}
\answer{
\begin{align*}
\textbf{Classificação:}\ SI
\end{align*}
Não possui solução.
}

\subquestion{
\[
\begin{cases}
x - 2y + z = 0 \\
2x + y - 2z = 1 \\
x + z = 2 \\
x - 3y - z = -3
\end{cases}
\]
}
\answer{
\begin{align*}
\textbf{Classificação:}\ SPD
\end{align*}
\[
(x,y,z)=(1,1,1)
\]
}

\subquestion{
\[
\begin{cases}
x - 2y + 3z + t = 10 \\
x - y - 2t = -9 \\
2x - 4y + 7z = 15 \\
y + z - t = 1
\end{cases}
\]
}
\answer{
\begin{align*}
\textbf{Classificação:}\ SPD
\end{align*}
\[
(x,y,z,t)=(1,2,3,4)
\]
}

\subquestion{
\[
\begin{cases}
x + 2y - z = 2 \\
2x + 5y + z = 8 \\
x - y - z = -1 \\
3x + y + z = 5
\end{cases}
\]
}
\answer{
\begin{align*}
\textbf{Classificação:}\ SPD
\end{align*}
\[
(x,y,z)=(1,1,1)
\]
}

\subquestion{
\[
\begin{cases}
x - 3y + 2z = 1 \\
x + y - z = 1 \\
3x - 5y + 3z = 3
\end{cases}
\]
}
\answer{
\begin{align*}
\textbf{Classificação:}\ SPI
\end{align*}
\[
x=1+\frac{z}{4}, \quad y=\frac{3z}{4}, \quad z\in\mathbb{R}
\]
}

\subquestion{
\[
\begin{cases}
x + 2y - z = 2 \\
2x - 3y + z = 0 \\
4x + y - z = 4
\end{cases}
\]
}
\answer{
\begin{align*}
\textbf{Classificação:}\ SPI
\end{align*}
\[
x=\frac{6}{7}+\frac{z}{7}, \quad y=\frac{4}{7}+\frac{3z}{7}, \quad z\in\mathbb{R}
\]
}

\subquestion{
\[
\begin{cases}
x - y + z = 4 \\
2x + y + 2z = 8 \\
x - 2y - z = 0 \\
3x - 5y - z = 6
\end{cases}
\]
}
\answer{
\begin{align*}
\textbf{Classificação:}\ SI
\end{align*}
Não possui solução.
}

\subquestion{
\[
\begin{cases}
x - 2y + z = 0 \\
2x - 3y + 2z = 1 \\
3x + y - z = 7 \\
x + 4y - 3z = 6
\end{cases}
\]
}
\answer{
\begin{align*}
\textbf{Classificação:}\ SPD
\end{align*}
\[
(x,y,z)=(2,1,0)
\]
}

\subquestion{
\[
\begin{cases}
x - 3y + z = -4 \\
x + 2y - 3z = 12 \\
2x - y - 2z = 8
\end{cases}
\]
}
\answer{
\begin{align*}
\textbf{Classificação:}\ SPI
\end{align*}
\[
x=\frac{28}{5}+\frac{7z}{5}, \quad y=\frac{16}{5}+\frac{4z}{5}, \quad z\in\mathbb{R}
\]
}

\subquestion{
\[
\begin{cases}
x + y - z = 1 \\
2x + 3y + z = 6 \\
-x + y + z = 1
\end{cases}
\]
}
\answer{
\begin{align*}
\textbf{Classificação:}\ SPD
\end{align*}
\[
(x,y,z)=(1,1,1)
\]
}

\subquestion{
\[
\begin{cases}
x + 2y = 2 \\
3x - y = 6 \\
4x + y = 3
\end{cases}
\]
}
\answer{
\begin{align*}
\textbf{Classificação:}\ SI
\end{align*}
Não possui solução.
}

\subquestion{
\[
\begin{cases}
x + 2y - z = 2 \\
x - y + 2z = 5 \\
3x - y - z = 2 \\
2x + 3y + z = 11
\end{cases}
\]
}
\answer{
\begin{align*}
\textbf{Classificação:}\ SI
\end{align*}
Não possui solução.
}

\subquestion{
\[
\begin{cases}
x + 2y - z = 3 \\
2x + y + z = 3 \\
x + y + 2z = 4
\end{cases}
\]
}
\answer{
\begin{align*}
\textbf{Classificação:}\ SPD
\end{align*}
\[
(x,y,z)=(0,2,1)
\]
}

\subquestion{
\[
\begin{cases}
x + 3y - z = 2 \\
x - y + 3z = 4 \\
3x + y + 5z = 10
\end{cases}
\]
}
\answer{
\begin{align*}
\textbf{Classificação:}\ SPI
\end{align*}
\[
x=\frac{7}{2}-2z,\quad y=z-\frac{1}{2},\quad z\in\mathbb{R}
\]
}

\subquestion{
\[
\begin{cases}
x - y - z = 4 \\
2x + y + 2z = 3 \\
2x - 5y - 6z = 0
\end{cases}
\]
}
\answer{
\begin{align*}
\textbf{Classificação:}\ SI
\end{align*}
Não possui solução.
}

\subquestion{
\[
\begin{cases}
x + y - z = 5 \\
2x + y - z = 20 \\
x + 2y - 4y = -5
\end{cases}
\]
}
\answer{
\begin{align*}
\textbf{Classificação:}\ SPD
\end{align*}
\[
(x,y,z)=(15,10,20)
\]
}

\subquestion{
\[
\begin{cases}
x + 3y + z = 0 \\
2x + 5y - z = 3 \\
x - y + 2z = -1 \\
2x + 9y - 2z = 4
\end{cases}
\]
}
\answer{
\begin{align*}
\textbf{Classificação:}\ SPD
\end{align*}
\[
(x,y,z)=(1,0,-1)
\]
}

\subquestion{
\[
\begin{cases}
x + y + z - t = 3 \\
2x + 3y - z + t = 1 \\
-x + y - z - 2t = -1 \\
x + 2y + 2z + t = 6
\end{cases}
\]
}
\answer{
\begin{align*}
\textbf{Classificação:}\ SPD
\end{align*}
\[
(x,y,z,t)=(0,1,2,0)
\]
}

\subquestion{
\[
\begin{cases}
x - y + z + t = 4 \\
2x - y - z = -3 \\
x - 2y + t = 1 \\
5x + z - t = 4
\end{cases}
\]
}
\answer{
\begin{align*}
\textbf{Classificação:}\ SPD
\end{align*}
\[
(x,y,z,t)=\left(0,-\frac{2}{3},\frac{11}{3},-\frac{1}{3}\right)
\]
}

\subquestion{
\[
\begin{cases}
x + y + z - t = 2 \\
2x - y - z - t = -1 \\
x - 2y - 2z = -3 \\
3x - 3y - 3z - t = -4
\end{cases}
\]
}
\answer{
\begin{align*}
\textbf{Classificação:}\ SPI
\end{align*}
\[
x=\frac{1}{3}+\frac{2t}{3},\quad y=\frac{5}{3}+\frac{t}{3}-z,\quad z,t\in\mathbb{R}
\]
}

\subquestion{
\[
\begin{cases}
x + y + z = 6 \\
2x + y + z = 8 \\
3x + 2y + 2z = 14 \\
3x + y + z = 10
\end{cases}
\]
}
\answer{
\begin{align*}
\textbf{Classificação:}\ SPI
\end{align*}
\[
x=2,\quad y=4-z,\quad z\in\mathbb{R}
\]
}

\newpage
\question{
Três pessoas foram a uma lanchonete. A primeira pagou R\$ 4,00 por dois
guaranás e um pastel; a segunda pagou R\$ 5,00 por um guaraná e dois pastéis; e
a terceira pagou R\$ 7,00 por dois guaranás e dois pastéis. Todos pagaram o
preço correto? Justifique sua resposta.
}
\answer{
Seja $g$ o preço do guaraná e $p$ o preço do pastel.

\[
\begin{cases}
2g + p = 4 \\
g + 2p = 5 \\
2g + 2p = 7
\end{cases}
\]

Das duas primeiras equações:
\[
(2g+p) - (g+2p) = 4 - 5 \Rightarrow g - p = -1 \Rightarrow g = p - 1
\]

Substituindo em $2g+p=4$:
\[
2(p-1)+p=4 \Rightarrow 3p=6 \Rightarrow p=2,\quad g=1
\]

Verificação na terceira:
\[
2g+2p = 2(1)+2(2)=6 \neq 7
\]

\textbf{Conclusão:} sistema inconsistente com o terceiro pagamento.
Logo, \textbf{não, a terceira pessoa não pagou o preço correto}.
}

\question{
Determine a função polinomial de grau dois qual que:
}
\subquestion{
$f(1)=0$, $f(3)=2$ e $f(4)=6$.
}
\answer{
Seja $f(x)=ax^2+bx+c$.

\[
\begin{cases}
a+b+c=0 \\
9a+3b+c=2 \\
16a+4b+c=6
\end{cases}
\]

Resolvendo:
\[
f(x)=x^2-3x+2
\]
}

\subquestion{
$p(1)=-1$, $p(2)=5$ e $p(-1)=7$.
}
\answer{
Seja $p(x)=ax^2+bx+c$.

\[
\begin{cases}
a+b+c=-1 \\
4a+2b+c=5 \\
a-b+c=7
\end{cases}
\]

Resolvendo:
\[
p(x)=3x^2-2x-2
\]
}

\newpage
\subquestion{
$h(-1)=2$, $h(1)=2$, $h(2)=2$, e $h(-2)=3$.
}
\answer{
Quatro condições com três coeficientes indicam verificação de consistência.

Se $h(x)=ax^2+bx+c$:

\[
\begin{cases}
a-b+c=2 \\
a+b+c=2 \\
4a+2b+c=2 \\
4a-2b+c=3
\end{cases}
\]

Das duas primeiras:
\[
b=0,\quad a+c=2
\]

Terceira:
\[
4a+c=2
\]

Subtraindo:
\[
(4a+c)-(a+c)=2-2 \Rightarrow 3a=0 \Rightarrow a=0
\]

Logo:
\[
c=2
\]

Quarta equação:
\[
4(0)-2(0)+2=2 \neq 3
\]

\textbf{Conclusão:} dados inconsistentes, não existe função quadrática que satisfaça todas as condições.
}

\question{
Suponha que uma cultura de peixes é iniciada com 100 kg de alevinos de um mês
de idade e que, enquanto a idade deles está entre 1 e 10 meses, sua massa pode
ser aproximada por um polinômio de grau 3. Se ao atingirem 2 meses sua massa
total é de 350 kg, aos 3 meses é de 1.020 kg e aos 4 meses é de 1.990 kg,
determine o polinômio de grau 3 em questão. Utilize este polinômio para estimar
a massa total dos peixes quando atingirem 10 meses de idade.
}
\answer{
Seja $m(t)=at^3+bt^2+ct+d$.

Dados:
\[
m(1)=100,\quad m(2)=350,\quad m(3)=1020,\quad m(4)=1990
\]

Sistema:
\[
\begin{cases}
a+b+c+d=100 \\
8a+4b+2c+d=350 \\
27a+9b+3c+d=1020 \\
64a+16b+4c+d=1990
\end{cases}
\]

Resolvendo:
\[
m(t)=115t^3-465t^2+650t-200
\]

Estimativa:
\[
m(10)=115(1000)-465(100)+650(10)-200
\]
\[
m(10)=115000-46500+6500-200=74700
\]

\textbf{Resposta final:} $m(10)=74.700$ kg.
}
